\section{Theory}

The paper does not need a long formal apparatus to motivate the experiments. Two simple statements are enough.

\begin{proposition}
A stationary key-side projection cannot represent emission history. In particular, if the same operator $k' = k + \alpha Pk$ is applied at every decoding step, then the operator cannot distinguish ``facet already used'' from ``facet still unused'' without additional state.
\end{proposition}

\noindent
This proposition is structural rather than statistical. The operator is time-invariant, so any step-specific notion of coverage has to come from the model's internal dynamics alone. The intervention itself contributes no state.

\begin{proposition}
The natural query-side counterpart of coverage is an operator that suppresses the projection of the current query onto already used facet directions. A practical history-free approximation is to suppress the projection onto the full ontology subspace.
\end{proposition}

\noindent
The first clause describes the ideal mechanism; the second clause matches the implementation evaluated in this paper. The experiment therefore tests whether even the coarse approximation already has the right sign on a two-tool benchmark.

\begin{theorem}
Let $\mathrm{LHS}$ denote the empirical attention-output perturbation and let $\mathrm{RHS}$ denote the local bound term from the project's perturbation analysis. On both Qwen and Llama, the observed median ratio $\mathrm{LHS}/\mathrm{RHS}$ is on the order of $10^{-8}$, so the diagnostic bound is numerically conservative on the evaluated interventions.
\end{theorem}

\noindent
The theorem is an empirical statement in this paper. Its role is not to prove end-to-end accuracy, but to show that the perturbation language used to discuss Q-side and K-side interventions is consistent with the measured scale of the induced output changes.

\begin{remark}
The theory here is intentionally narrower than the older drafts in this repository. We do not claim that the current implementation realizes exact step-adaptive facet tracking, and we do not use unverified compression or predictor results as part of the main theorem chain.
\end{remark}
